\documentclass[12pt, letterpaper]{article} 

\usepackage{amsmath, amssymb, amsthm, enumerate}

\title{MATH 102 - Applied Linear Algebra}
\author{Dani Nguyen}
\date{Fall 2025}

\begin{document}
\maketitle
\tableofcontents

\section{Invertible Matrix Theorem}
\subsection{Definition}
Let A be a square n x n matrix. Then the following statements are
equivalent. That is, for a given A, the statements are either all 
true or all false.
\begin{enumerate}
    \item A is invertible.
    \item A is row equivalent to I.
    \item A has n pivot positions.
    \item The equation A\textbf{x} = \textbf{0} has only the trivial solution.
    \item The columns of A are linearly independent.
    \item The linear transformation \textbf{x} $\mapsto$ A\textbf{x} is one-to-one.
    \item The equation A\textbf{x} = \textbf{b} has at least one solution for each \textbf{b} in $\mathbb{R}^n$.
    \item The columns of A span $\mathbb{R}^n$
    \item The linear transformation \textbf{x} $\mapsto$ A\textbf{x} maps $\mathbb{R}^n$ onto $\mathbb{R}^n$.
    \item There is an n x n matrix C such that CA = I.
    \item There is an n x n matrix D such that AD = I.
    \item A$^T$ is invertible.
    \item The columns of A form a basis of $\mathbb{R}^n$.
    \item Col A = $\mathbb{R}^n$.
    \item dim Col A = n.
    \item rank A = n
    \item Nul A = {0}
    \item dim Nul A = 0
    \item The number 0 is \emph{not} an eigenvalue of A.
    \item det A $\neq$ 0.
    \item (Col A)$^\bot$ = \textbf{0}
    \item (Nul A)$^\bot$ = $\mathbb{R}^n$
    \item Row A = $\mathbb{R}^n$
\end{enumerate}

\subsection{Exercise}
\textbf{Thm 9.1.16} Let A $\in M_n$ and let $\lambda \in \mathbb{C}$. Then the following statements are equivalent:
\begin{enumerate}[(a)]
    \item $\lambda$ is an eigenvalue of A.
    \item A\textbf{x} = $\lambda$\textbf{x} for some nonzero \textbf{x} $\in \mathbb{C}^n$.
    \item (A - $\lambda$I)\textbf{x} = \textbf{0} has a nontrivial solution, that is, nullity(A - $\lambda I$) $>$ 0.
    \item rank(A - $\lambda$I) $<$ n.
    \item A - $\lambda$I is not invertible.
    \item A$^\top$ - $\lambda$I is not invertible.
    \item $\lambda$ is an eigenvalue of A$^\top$.
\end{enumerate}
\textbf{Proof}
\begin{enumerate}[(a)]
    \item $\Leftrightarrow$ (b) 
    
    By definition 9.1.1, if $A\mathbf{x}=\lambda\mathbf{x}$ and $\mathbf{x}\neq\mathbf{0}$ then $(\lambda,\mathbf{x})$
    is an eigenpair of A, meaning $\lambda$ is an eigenvalue of A.
    \item $\Leftrightarrow$ (c)
    
    These are restatements of one another.
    \item $\Leftrightarrow$ (d)
    
    By rank nullity theorem, if nullity(A - $\lambda I$) $>$ 0, then rank(A - $\lambda I$) $<$ n.
    \item $\Leftrightarrow$ (e) 
    
    By property 1 and 16 of IMT, if rank(A - $\lambda I$) $<$ n then (A - $\lambda I$) is not invertible.
    \item $\Leftrightarrow$ (f)
    
    Suppose we have $A - \lambda I = $ 
    $\begin{bmatrix}
        a - \lambda I & b  \\
        c & d - \lambda I  \\
    \end{bmatrix}$, 
    and suppose it is not invertible, its determinant is $(a - \lambda I)(d - \lambda I) - bc = 0$. Consider the case of $A^\top - \lambda I$, we have 
    $\begin{bmatrix}
        a - \lambda I & c  \\
        b & d - \lambda I  \\
    \end{bmatrix}$, 
    whose determinant is $(a - \lambda I)(d - \lambda I) - cb$. And by commutativity of multiplication, $bc = cb$, so the determinant is also 0, meaning $A^\top - \lambda I$ is also not invertible.

    \item $\Leftrightarrow$ (g)
    
    Since (a) is equivalent to (e), the same is true for $A^\top$.
\end{enumerate}

\section{Annihilator}
\subsection{Some Roots of an Annihilator are Eigenvalues}
While all the eigenvalues of a matrix are roots of each of its annihilators,
the roots of an annihilator are not necessarily eigenvalues. \\

\noindent This raises a question:

Must \underline{some} of the roots of an annihilator of a matrix be eigenvalues
of the matrix? The answer is yes. \\

\noindent Let $p(z) = a_0 + a_1z + a_2z^2 + \ldots + a_kz^k$ be an annihilator of a matrix $A \in \mathbb{M}_n$. \\
Factor $p$ as follows: $p(z) = a_k(z-\lambda _1)(z-\lambda _2)\ldots(z-\lambda _k), a_k \ne 0$ \\
Since $p$ annihilates $A$, we have
\[
0=p(A)=a_k(A-\lambda_1I)(A-\lambda_2I)\ldots(A-\lambda_kI)
\]
Since the zero matrix is not invertible, one of the factors $A-\lambda_iI$ is not invertible. 
This ensures that $\lambda_i$ is an eigenvalue of $A$.

\subsection{Each Matrix has an Eigenvector}
\subsubsection{Proof 1}
$\lambda$ is an eigenvector of $A \in \mathbb{M}_n$ iff det($A-\lambda I$) = 0. \\
$p(\lambda) = \text{det} (A-\lambda I)$ is a polynomial of degree n (the characteristic polynomial of A). \\
Therefore, $\lambda$ is an eigenvalue of A iff $\lambda$ is a root of p (i.e. $p$) \\
By the Fundamental Theorem of Algebra (FTA), each polynomial over the complex has a root. \\
Hence, each matrix $A \in \mathbb{M}_n$ has an eigenvalue.
\subsubsection{Proof 2}
We've shown that each matrix $A \in \mathbb{M}_n$ has an annihilator $p$. \\
By the FTA, $p$ must have a root. \\
In 2.1, we showed that at least one root of $p$ is an eigenvalue of $A$. \\
Hence, each matrix $A \in \mathbb{M}_n$ has an eigenvalue. 

\section{Eigenvalues}
\subsection{Commuting Matrices Share An Eigenvector}
\subsubsection{2-Commute}
\textbf{Thm:} Let $A_1, A_2 \in \mathbb{M}_n$. Let $\lambda$ be an eigenvalue of $A_1$. 

Suppose $A_1A_2=A_2A_1$, then $A_1,A_2$ share an eigenvector in $\mathcal{E}_\lambda (A_1)$ \\

\noindent \textbf{Computation:} 
\begin{enumerate}
    \item[0.] Choose $A_1$ or $A_2$. Say, $A_1$
    \item Compute an eigenvalue of $A_1$. Say, $\lambda$, and form its eigenspace $\mathcal{E}_\lambda(A_1)$
    \item Select a basis $\ \vec x_1,\vec x_2,...,\vec x_k\ $ for $\mathcal{E}_\lambda(A_1)$ and form a matrix $X=[\vec x_1\; \vec x_2\; \ldots\; \vec x_k]$
    \item Solve $A_2X=XC$ for $C$
    \item Find an eigenpair of $C:(\mu,u)$
    \item Compute $v=Xu$
\end{enumerate}
\newpage
\noindent \textbf{Proof:} 
\begin{align*}
    Ax_1=\lambda x_1,\ \ Ax_2=\lambda x_2,\ ...\; ,\ Ax_k=\lambda x_k \\
    A\vec x_i=\lambda \vec x_i\qquad\text{for}\ i=1,2,...,k 
\end{align*}
\begin{align*}
    A_1x_i &= \lambda x_i \\
    A_2(A_1x_i) &= A_2(\lambda x_i) \\
    A_1(A_2x_i) &= \lambda (A_2x_i) \tag{commutativity}\\
\end{align*}
\begin{enumerate}
    \item[Case 1:] $A_2x_i=0$ \\
                    $\Rightarrow A_2x_i \in \mathcal{E}_\lambda(A_1)$ 
    \item[Case 2:] $A_2x_i\ne0$ \\
                    $\Rightarrow A_2x_i$ is an eigenvector of $A_1$ and 
                    $A_2x_i \in \mathcal{E}_\lambda(A_1)$
\end{enumerate}
\begin{align*}
    A_2x_1=Xc_1,\ A_2x_2&=Xc_2,\ ...\; ,\ A_2x_k=Xc_k \\
    A_2[x_1\ x_2\ ...\ x_k]=[A_2x_1\ A_2x_2\ ...\ A_2x_k]&=[Xc_1\ Xc_2\ ...\ Xc_k]=X[c_1\ c_2\ ...\ c_k] \\
    A_2X&=XC 
\end{align*}
\begin{align*}    
    \text{Let } (\mu,u) \text{ be an eigenpair of } C & \qquad C\vec u=\mu \vec u \qquad u\ne 0 \\
    \text{Say } v&=Xu \ne 0 \\
    Cu &= \mu u \\
    XCu &= X\mu u \\
    A_2Xu &= \mu Xu \\
    A_2v &= \mu v \\
    \therefore v=Xu \in \mathcal{E}_\lambda(A_1) &\text{ is an eigenvector of } A_2
\end{align*}

\end{document}